\section{Mathematical Formulation}
\label{sec:formulation}

% =============================================================================
% 2.1  Geometry and coordinate system
% =============================================================================
\subsection{Equivalent-sphere shell model}
\label{sec:geometry}

We model the watermelon as an oblate spheroid with semi-major axis $a$
(equatorial radius), semi-minor axis $c$ (polar radius), and uniform
rind thickness $h$.  The aspect ratio is defined as
$\zeta = c/a$,
where $\zeta < 1$ corresponds to the characteristic oblate shape of most
commercial watermelon cultivars ($\zeta \approx 0.5$--$0.7$ for
elongated varieties; $\zeta \approx 0.8$--$0.9$ for round cultivars).
The equivalent-volume sphere has radius
\begin{equation}
  R_{\mathrm{eq}} = \left( a^2 c \right)^{1/3}.
  \label{eq:Req}
\end{equation}

\noindent
The eigenfrequency model developed below is an \emph{equivalent-sphere
approximation}: all geometry enters through~$R_{\mathrm{eq}}$, and the
frequency equation is that of a spherical shell of
radius~$R_{\mathrm{eq}}$.  The model does not include explicit
curvature corrections for oblate deformation (i.e.\ there are no
$\zeta$-dependent geometric factors beyond those captured
by~$R_{\mathrm{eq}}$).  Consequently, two watermelons with the same
$R_{\mathrm{eq}}$ but different aspect ratios will be assigned the
same predicted frequency.

This approximation is sufficient for the inversion application developed
in~\S\ref{sec:inversion}: the goal is to recover~$E_{\mathrm{rind}}$
from a measured~$f_2$ at \emph{known} geometry.  Since the fruit
dimensions~$a$, $c$, and~$h$ are measured independently,
$R_{\mathrm{eq}}$ is determined, and the inversion
$E = f_2^2 / \alpha$ is well-posed regardless of whether the forward
model captures shape effects beyond~$R_{\mathrm{eq}}$.  Nevertheless,
the absence of explicit curvature corrections means that the model
cannot predict how $f_2$ would differ between a sphere and an oblate
spheroid of the same equivalent volume.  We discuss this limitation
further in~\S\ref{sec:limitations} and identify a Ritz variational
extension as future work.

At fixed equatorial radius~$a$, varying~$\zeta$ from~0.5 to~0.95
changes~$R_{\mathrm{eq}}$ (because $c = \zeta a$ changes the
equivalent volume) and thereby shifts~$f_2$ by
approximately~\SI{30}{\percent}.  This is a \emph{size} effect
mediated by~$R_{\mathrm{eq}}$, not a \emph{shape} effect arising
from curvature anisotropy.  The distinction matters: the model
requires separate specification of both~$a$ and~$c$ only because
they jointly determine~$R_{\mathrm{eq}}$, not because the frequency
equation contains independent $a$- and $c$-dependent terms.

The rind is treated as a thin, isotropic, viscoelastic shell with Young's
modulus $E_{\mathrm{rind}}$, Poisson's ratio $\nu$, density
$\rho_w$, and structural loss factor $\eta$.  The enclosed flesh is
modelled as an inviscid, nearly incompressible fluid with density
$\rho_f$ and bulk modulus $K_f$.  The canonical parameter values used
throughout this study are given in Table~\ref{tab:canonical}.

\begin{table}[t]
  \centering
  \caption{Canonical parameters for the Crimson Sweet watermelon at
    optimal ripeness.  Values are drawn from Sadrnia et
    al.~\cite{Sadrnia2008}, De Belie et al.~\cite{DeBelieetal2000},
    and established tissue mechanics compilations.}
  \label{tab:canonical}
  \begin{tabular}{llcc}
    \toprule
    Parameter & Symbol & Value & Unit \\
    \midrule
    Semi-major axis    & $a$                 & 0.158 & m \\
    Semi-minor axis    & $c$                 & 0.123 & m \\
    Aspect ratio       & $\zeta = c/a$       & 0.78  & --- \\
    Equivalent radius  & $R_{\mathrm{eq}}$   & 0.145 & m \\
    Rind thickness     & $h$                 & 15    & mm \\
    Rind modulus (ripe) & $E_{\mathrm{rind}}$ & 50   & MPa \\
    Poisson's ratio    & $\nu$               & 0.40  & --- \\
    Rind density       & $\rho_w$            & 1050  & kg/m$^3$ \\
    Flesh density      & $\rho_f$            & 950   & kg/m$^3$ \\
    Flesh bulk modulus  & $K_f$              & 2.2   & GPa \\
    Turgor pressure     & $P_{\mathrm{int}}$ & 500   & Pa \\
    Loss tangent        & $\eta$             & 0.15  & --- \\
    \bottomrule
  \end{tabular}
\end{table}

% Schematic figure (placeholder — no diagram generated yet)
% \begin{figure}[t]
%   \centering
%   \includegraphics[width=0.6\textwidth]{fig_watermelon_schematic}
%   \caption{Geometry of the oblate spheroidal shell model.  The
%     watermelon is characterised by semi-major axis~$a$, semi-minor
%     axis~$c$, rind thickness~$h$, and flesh properties
%     ($\rho_f$, $K_f$).  The coordinate system and modal
%     decomposition are indicated.}
%   \label{fig:schematic}
% \end{figure}

% =============================================================================
% 2.2  Eigenfrequency derivation
% =============================================================================
\subsection{Eigenfrequency of the flexural mode}
\label{sec:eigenfreq}

For a thin spherical shell of radius~$R_{\mathrm{eq}}$ filled with
fluid, the natural frequency of the $n$th flexural mode is determined
by the balance of bending, membrane, and pre-stress stiffness against
the effective inertia of the shell--fluid system.  The effective mass
per unit area, including fluid loading, is
\begin{equation}
  \rho_{\mathrm{eff}} = \rho_w h + \frac{\rho_f\, R_{\mathrm{eq}}}{n},
  \label{eq:rho_eff}
\end{equation}
where the second term represents the added mass of the entrained fluid
for mode~$n$ \cite{Junger1986, lamb1882}.  Note that
$\rho_{\mathrm{eff}}$ has units of
\si{\kilogram\per\metre\squared} (mass per unit area).  To construct a
dimensionless group (\S\ref{sec:Pi_ripe}), we instead use
the volumetric effective density
\begin{equation}
  \tilde{\rho}_{\mathrm{eff}} = \frac{\rho_w\, h}{R_{\mathrm{eq}}} + \rho_f,
  \label{eq:rho_eff_vol}
\end{equation}
which has units of \si{\kilogram\per\metre\cubed}.  All occurrences of
$\rho_{\mathrm{eff}}$ in the dimensionless analysis
(Eqs.~\ref{eq:Pi_ripe}--\ref{eq:Pi_ripe_universal}) refer to this
volumetric form.

The dominant tap-tone mode is the lowest flexural mode, $n = 2$.  The
modal stiffness decomposes as
\begin{equation}
  K_n = K_{\mathrm{bend}} + K_{\mathrm{memb}} + K_{\mathrm{prestress}},
  \label{eq:stiffness_decomp}
\end{equation}
where the bending stiffness scales as
\begin{equation}
  K_{\mathrm{bend}} = \frac{n(n-1)(n+2)^2\, h^3}
    {12(1-\nu^2)\, R_{\mathrm{eq}}^4}\, E_{\mathrm{rind}},
  \label{eq:K_bend}
\end{equation}
the membrane stiffness scales as
\begin{equation}
  K_{\mathrm{memb}} = \frac{h}{R_{\mathrm{eq}}^2}\,
    \frac{n^2 + n - 2 + 2\nu}{n^2 + n + 1 - \nu}\,
    E_{\mathrm{rind}},
  \label{eq:K_memb}
\end{equation}
and the pre-stress stiffness (from internal turgor pressure~$P_{\mathrm{int}}$) is
\begin{equation}
  K_{\mathrm{prestress}} = \frac{P_{\mathrm{int}}}{R_{\mathrm{eq}}}
    \, (n-1)(n+2).
  \label{eq:K_prestress}
\end{equation}

The squared natural frequency is then
\begin{equation}
  f_n^2 = \frac{K_n}{4\pi^2\, \rho_{\mathrm{eff}}}.
  \label{eq:f2_squared}
\end{equation}

\noindent
Equations~\eqref{eq:K_bend}--\eqref{eq:K_prestress} are the standard
Lamb--Junger expressions for a \emph{spherical} shell
\cite{lamb1882, Junger1986, Leissa1973, Soedel2004}.  All geometry
enters through~$R_{\mathrm{eq}}$; the equations contain no explicit
dependence on the aspect ratio~$\zeta$.  The model is therefore an
equivalent-sphere approximation: it assigns to each oblate watermelon
the frequency of a sphere with the same enclosed volume.

\noindent A key observation: $f_2^2$ is \emph{linear} in
$E_{\mathrm{rind}}$ when turgor pressure is negligible (the residual
turgor contribution is less than \SI{0.5}{\percent} of the total
stiffness at canonical parameters).  This linearity
is the foundation of the closed-form inversion developed in
\S\ref{sec:inversion}.

\paragraph{Operational identification of the $n = 2$ mode.}
In the equivalent-sphere model, the target mode is the lowest flexural
mode ($n = 2$), which produces a quadrupolar deformation pattern---the
shell elongates along one axis while contracting along the
perpendicular axis.  Experimentally, this mode would be identified as
the lowest-frequency spectral peak in the impulse response of the
tapped fruit (typically obtained via FFT of a microphone or
accelerometer signal).  For watermelons in the ripe range, this
corresponds to a frequency of approximately \SIrange{80}{110}{\hertz}
(Table~\ref{tab:forward_results}).  Higher flexural modes ($n = 3, 4$)
produce additional spectral peaks at higher frequencies but carry less
energy in a knuckle-tap excitation.  Confirmation of mode identity
would require multi-point vibration measurement (e.g.\ scanning laser
Doppler vibrometry), which is beyond the scope of the present
analytical study.

For the canonical ripe watermelon ($n = 2$, $h = \SI{15}{\milli\metre}$,
$R_{\mathrm{eq}} = \SI{0.145}{\metre}$), the bending and membrane
terms both scale linearly with~$E_{\mathrm{rind}}$, whereas the
pre-stress term is independent of~$E$.  This separation underpins the
linearity $f_2^2 = \alpha E_{\mathrm{rind}}$, where the proportionality
constant is
\begin{equation}
  \alpha = \frac{K_{\mathrm{bend}}^{(0)} + K_{\mathrm{memb}}^{(0)}}
    {4\pi^2\, \rho_{\mathrm{eff}}}
    = \SI{161.5}{\hertz\squared\per\mega\pascal}
  \label{eq:alpha_value}
\end{equation}
for the canonical geometry, where the superscript~$(0)$ denotes the
geometric pre-factors of~$E_{\mathrm{rind}}$ in
Eqs.~\eqref{eq:K_bend}--\eqref{eq:K_memb}.

% =============================================================================
% 2.3  Closed-form inversion
% =============================================================================
\subsection{Inversion: from tap-tone to rind modulus}
\label{sec:inversion}

Rearranging Eq.~\eqref{eq:f2_squared} for $E_{\mathrm{rind}}$ gives
the explicit inversion
\begin{equation}
  \boxed{
    E_{\mathrm{rind}} = \frac{\omega^2\, m_{\mathrm{eff}}
      - K_{\mathrm{prestress}}}
      {K_{\mathrm{bend}}^{(0)} + K_{\mathrm{memb}}^{(0)}},
  }
  \label{eq:inversion}
\end{equation}
where $\omega = 2\pi f_2$, $m_{\mathrm{eff}} = \rho_w h + \rho_f
R_{\mathrm{eq}} / n$ is the effective mass per unit area including
fluid loading, and $K_{\mathrm{bend}}^{(0)}$ and
$K_{\mathrm{memb}}^{(0)}$ are the geometric pre-factors of
$E_{\mathrm{rind}}$ in Eqs.~\eqref{eq:K_bend}--\eqref{eq:K_memb}.
Thus a single measurement of the tap-tone frequency~$f_2$, combined
with knowledge (or estimation) of the fruit geometry and flesh density,
yields the rind elastic modulus directly.  No iterative solver or
optimisation is required.

Equation~\eqref{eq:inversion} is algebraically exact within the
equivalent-sphere model (which itself is an approximation to the true
oblate geometry): the
round-trip error $|E_{\mathrm{inv}} - E_{\mathrm{true}}| /
E_{\mathrm{true}}$ is less than \SI{0.01}{\percent} for all ripeness
stages tested (\S\ref{sec:forward}).

The condition number of this inversion with respect to~$f_2^2$ is unity,
since the mapping $f_2^2 \mapsto E_{\mathrm{rind}}$ is linear.  The
practical condition number with respect to the \emph{measured
frequency}~$f_2$ is approximately~2, because $E \propto f_2^2$ implies
$\delta E / E \approx 2\, \delta f / f$.  This~$2\times$ noise
amplification is inherent to any square-law inversion and is confirmed
numerically in the Monte Carlo analysis of \S\ref{sec:geom_uncertainty}.

% =============================================================================
% 2.4  Dimensional analysis and the ripeness parameter
% =============================================================================
\subsection{Dimensionless geometric invariant}
\label{sec:Pi_ripe}

To collapse predictions across cultivars, we define the dimensionless
group
\begin{equation}
  \Pi_{\mathrm{ripe}} = f_2 \, R_{\mathrm{eq}}
    \sqrt{\frac{\tilde{\rho}_{\mathrm{eff}}}{E_{\mathrm{rind}}}},
  \label{eq:Pi_ripe}
\end{equation}
where $\tilde{\rho}_{\mathrm{eff}} = \rho_w h / R_{\mathrm{eq}} +
\rho_f$ is the volumetric effective density
(Eq.~\ref{eq:rho_eff_vol}; units \si{\kilogram\per\metre\cubed}),
ensuring that $\Pi_{\mathrm{ripe}}$ is dimensionless.
From Eq.~\eqref{eq:f2_squared}, this parameter evaluates to
\begin{equation}
  \Pi_{\mathrm{ripe}} = \frac{1}{2\pi}
    \sqrt{\frac{(K_{\mathrm{bend}}^{(0)} + K_{\mathrm{memb}}^{(0)})
    \cdot R_{\mathrm{eq}}^2 \cdot \tilde{\rho}_{\mathrm{eff}}}
    {\rho_{\mathrm{eff}}}}
  \label{eq:Pi_ripe_universal}
\end{equation}
for fixed geometry (neglecting the pre-stress contribution).

A key observation is that $E_{\mathrm{rind}}$ cancels from
$\Pi_{\mathrm{ripe}}$ \emph{by construction}: since $f_2 \propto
\sqrt{E_{\mathrm{rind}}}$, the modulus appears in both the numerator
($f_2$) and denominator ($\sqrt{E_{\mathrm{rind}}}$) and divides out.
This cancellation is exact only when the geometry---including the rind
thickness~$h$---is held fixed while $E_{\mathrm{rind}}$ varies.
Consequently, $\Pi_{\mathrm{ripe}}$ is \emph{not} invariant across
ripeness stages in which both~$E$ and~$h$ change simultaneously
(as occurs during biological ripening;
Table~\ref{tab:forward_results}).
Rather, $\Pi_{\mathrm{ripe}}$ is a
\emph{geometric invariant} that characterises the shell--fluid system
at a given geometry, independently of the elastic modulus.

The practical utility of $\Pi_{\mathrm{ripe}}$ lies in
cultivar-independent calibration: it collapses different cultivar
\emph{geometries} onto a single calibration constant.  Once its value is
known for a given geometry, any measured $f_2$ can be mapped to
$E_{\mathrm{rind}}$ via rearrangement of Eq.~\eqref{eq:Pi_ripe},
without detailed knowledge of the flesh properties or internal pressure.
The universality is across cultivars (different geometries collapsing
onto a common constant), not across ripeness states---the latter would
require that the geometry remain fixed during ripening, which is only
approximately true for $a$ and~$c$ but not for~$h$.

For the four cultivars considered in this study (Crimson Sweet, Sugar
Baby, Charleston Gray, and Seedless Tri-X~313),
$\Pi_{\mathrm{ripe}} = 0.062 \pm 0.002$ (CoV~$= \SI{3.3}{\percent}$;
see Table~\ref{tab:Pi_ripe} in \S\ref{sec:universal_curve}).  This
tight collapse confirms that the dimensionless group effectively
normalises out the dominant cultivar-specific geometric variation.
